% Appendix A: the inventory behind section 1.1 (article-kit WRITING.md section 4). It adds detail
% to what 1.1 says and holds no caveat that 1.1 does not state. The table is read from the
% boundary harness's pins (Formalization/CIAxiomGuard.lean, the #guard_msgs blocks of the
% headline declarations of linkage.toml's [boundary] list), each axiom mapped to its cited fact
% through blueprint/AXIOMS.md's **Lean:** lines. The map of the fourteen axioms is
% blueprint/trust-boundary.txt's list of admitted names. The per-statement table's rows,
% appendix-axioms-table.tex, are generated from `#print axioms` over every declaration a printed
% statement names, by scripts/theory_axioms_table.py (its docstring has the commands; run where the
% Lean project is built). Regenerate them after any change to a statement's declarations.
\section{Verification}
\label{app:verification}

The formal development is in Lean~4 (\texttt{v4.31.0}) on Mathlib at its tag \texttt{v4.31.0}
(commit \texttt{fabf563}). The script \texttt{scripts/theory\_axioms\_table.py} of the
accompanying repository regenerates Table~\ref{tab:statements} from Lean's \texttt{\#print
axioms}, and its header gives the commands.
% The export carries the scripts the paper names by path through the module table's extra_files
% (linkage.toml), since by itself it copies only scripts/make-fig-*.py.
The development imports the formal
developments of the author's earlier papers at their releases: the shared core first released with
the line paper \progpaper{V-line v0.1}, \texttt{ScaleSpaceCore}, here at its later tag
\texttt{v0.5.0}, which declares no axiom, and
the export \texttt{AffineHemigroup} of the affine paper \progpaper{affine v0.1}, at \texttt{v0.1}.

Table~\ref{tab:axioms} maps the fourteen axioms of the development to the nine admitted cited
facts of Section~\ref{sec:trust}. Six of them carry the namespace \texttt{AffineHemigroup}, the
affine paper's export, and eight the namespace \texttt{Pyramids}, this development. A name
\texttt{cited\_A}$n$ refers to entry $n$ of the project's ledger of cited facts, which the paper
prints as a numbered theorem of Section~\ref{ch:cited}: \texttt{cited\_A48} is
Theorem~\ref{thm:cited-A48}, and the table gives the theorem of every axiom. The supplement's
Section~\ref{supp:axioms} adds the definitions the axioms name.

\begin{table}[ht]
\centering
\small
\begin{tabular}{ll}
\toprule
Axiom & Cited fact \\
\midrule
\texttt{levy\_khintchine\_exists}, \texttt{\_unique}, \texttt{\_converse} & \ledger{A39} \\
\texttt{khintchine\_null\_array} & \ledger{A41} \\
\texttt{selfDecomposable\_infinitelyDivisible} & \ledger{A43}, first clause \\
\texttt{cited\_A43\_residual\_infinitelyDivisible} & \ledger{A43}, the residual \\
\texttt{cited\_A43\_semi\_infinitelyDivisible} & \ledger{A43}, the semi-selfdecomposable law \\
\texttt{sato\_polar\_criterion} & \ledger{A45} \\
\texttt{cited\_A46} & \ledger{A46} \\
\texttt{cited\_A48} & \ledger{A48} \\
\texttt{cited\_A49} & \ledger{A49} \\
\texttt{cited\_A22\_polya}, \texttt{cited\_A22\_variation} & \ledger{A22} \\
\texttt{cited\_A24} & \ledger{A24} \\
\bottomrule
\end{tabular}
\caption{The fourteen axioms of the formal development and the cited facts they state. Six are
axioms of \texttt{AffineHemigroup}: the four of the first two rows,
\texttt{selfDecomposable\_infinitelyDivisible} and \texttt{sato\_polar\_criterion}. The other eight
are axioms of \texttt{Pyramids}.}
\label{tab:axioms}
\end{table}

Table~\ref{tab:headline} lists the headline declarations, the statement each one formalizes and the
cited facts that Lean's \texttt{\#print axioms} reports for it beyond its three logical axioms, read
through Table~\ref{tab:axioms}. They are formal counterparts of main results of
Sections~\ref{ch:pyramid}, \ref{ch:matern} and~\ref{ch:price} and of two statements of the
supplement, and their lists
of axioms are fixed in the build: a Lean file prints the axioms of each and compares them with a
stored copy, so that the build fails when a list changes.

\begin{table}[!htbp]
\centering
\small
\begin{tabular}{>{\raggedright\arraybackslash}p{0.36\textwidth}l>{\raggedright\arraybackslash}p{0.27\textwidth}}
\toprule
Declaration & Statement & Cited facts \\
\midrule
\texttt{lattice\_strata\_1}, \texttt{\_2}, \texttt{\_3} & Theorem~\ref{thm:lattice-strata}(1)--(3) & \ledger{A39}, \ledger{A41} \\
\texttt{lattice\_strata\_strict} & Theorem~\ref{thm:lattice-strata-strict} & \ledger{A39}, \ledger{A41}, \ledger{A45} \\
\texttt{what\_does\_not\_lift\_chain\_strict} & Proposition~\ref{prop:what-does-not-lift-chain} & \ledger{A39}, \ledger{A43} \\
\texttt{ladder\_knot\_kernel\_1}, \texttt{\_2} & Proposition~\ref{prop:ladder-knot-kernel}(1), (2) & \ledger{A39}, \ledger{A41} \\
\texttt{ladder\_knot\_kernel\_2\_moments} & Proposition~\ref{prop:ladder-knot-kernel}(2), moments & \ledger{A49} \\
\texttt{ladder\_knot\_kernel\_2\_unique\_iff}, \texttt{\_rot} & Proposition~\ref{prop:ladder-knot-kernel}(2), uniqueness & \ledger{A39}, \ledger{A41} \\
\texttt{no\_bounded\_step} & Proposition~\ref{prop:no-bounded-step} & \ledger{A41}, \ledger{A46} \\
\texttt{matern\_pyramid\_1\_family} & Theorem~\ref{thm:matern-pyramid}(1) & \ledger{A39}, \ledger{A43}, \ledger{A45} \\
\texttt{creation\_plane\_1} & Proposition~\ref{prop:creation-plane}(1) & none \\
\texttt{stage\_lemma\_3} & Proposition~\ref{prop:stage-lemma}(3) & none \\
\texttt{selection\_stratum\_two\_2}, \texttt{\_converse} & Theorem~\ref{thm:selection-stratum-two}(2) & none \\
\texttt{grid\_pyramid\_not\_isotropic\_2} & Proposition~\ref{prop:grid-pyramid-not-isotropic}(2) & \ledger{A39}, \ledger{A41}, \ledger{A48} \\
\texttt{exact\_masks\_1} & Proposition~\ref{prop:exact-masks}(1), the criterion & \ledger{A39}, \ledger{A41}, \ledger{A49} \\
\texttt{exact\_masks\_2} & Proposition~\ref{prop:exact-masks}(2) & \ledger{A39}, \ledger{A41}, \ledger{A48} \\
\texttt{discrete\_helmholtz\_kernel}, \texttt{\_cascade} & Proposition~\ref{prop:discrete-helmholtz}(1), (2) & none \\
\texttt{lattice\_line\_1}, \texttt{\_2}, \texttt{\_3} & Theorem~\ref{thm:lattice-line} & \ledger{A39}, \ledger{A41} \\
\texttt{lattice\_line\_strict} & Theorem~\ref{cor:lattice-line-strict} & \ledger{A39}, \ledger{A41} \\
\bottomrule
\end{tabular}
\caption{The headline declarations (namespace \texttt{Pyramids}) and the cited facts Lean reports
for them. The last two rows are statements of the supplement.}
\label{tab:headline}
\end{table}
\clearpage

Table~\ref{tab:statements} gives, for every numbered statement of the paper and of the supplement,
the cited facts that Lean's \texttt{\#print axioms} reports for its declarations, read through
Table~\ref{tab:axioms}. Each row is one section and one set of cited facts, and it lists the
statements whose declarations read exactly that set; ``none'' means Lean's three logical axioms
only. A definition is listed with the statements of its section. It has nothing to prove: its row
says only that its formal counterpart reads no axiom, and that counterpart was compared with the
printed definition. Section~\ref{sec:trust} names what is proved in print only: one step of
Proposition~\ref{prop:grid-pyramids-infinitely-divisible}(4), whose row lists the cited facts of
its declarations; clause (3) of Proposition~\ref{prop:grid-pyramid-singular-kernels}; and
Propositions~\ref{prop:expanding-grid-pyramid-not-isotropic}
and~\ref{prop:band-limited-refinable-mask}. Seven cited facts have no row:
Theorems~\ref{thm:cited-A30}, \ref{thm:cited-A13}, \ref{thm:cited-A15}, \ref{thm:cited-A19},
\ref{thm:cited-A20}, \ref{thm:cited-A55} and~\ref{thm:cited-A57} are stated for the reader and read
by no formal proof. Their statements are typed, without proof, in a skeleton library that the formal
development does not import. Propositions~\ref{prop:expanding-grid-pyramid-not-isotropic},
\ref{prop:grid-pyramid-singular-kernels} and~\ref{prop:band-limited-refinable-mask} have no row
either, since each has a clause without a formal proof. Those clauses are typed in the same skeleton
library without a proof, and no formal proof reads them. The first two clauses of
Proposition~\ref{prop:grid-pyramid-singular-kernels} are proved, by the declarations
\texttt{grid\_pyramid\_singular\_kernels\_1} and \texttt{\_2}, for which Lean reports
\ledger{A39} and \ledger{A41}: the printed proof reads \ledger{A40}, which the development proves
from \ledger{A39}.

A row lists the cited facts of a statement's own formal proof. Some formal statements take as a
hypothesis an object that an earlier statement builds. Theorem~\ref{thm:lattice-strata} takes the
data $(a(t)I, \nu_t)$ of an isotropic family, and Propositions~\ref{prop:matern-generator-octave}
and~\ref{prop:creation-plane} and Theorem~\ref{thm:matern-pyramid-line-noncreation} take a family
with the transforms of the \Matern{} pyramid. Theorem~\ref{thm:knot-atom-cascades},
Propositions~\ref{prop:strata-cylinder-coordinates} and~\ref{prop:phase-entropy-scale-unit} and
Lemma~\ref{lem:phase-disintegration-exists} take the pair of Theorem~\ref{thm:surface-measure}.
Such a statement is quantified over every object with the defining property, and its row lists only
its own cited facts. The facts that build the object are in the row of the statement that builds
it. A cited fact that the development proves with no cited fact, such as
Theorem~\ref{thm:cited-A51}, adds nothing to a row, and Lemma~\ref{lem:bessel-k}, whose printed
proof reads it, has the row ``none''.

{\small
\begin{longtable}{l>{\raggedright\arraybackslash}p{0.24\textwidth}>{\raggedright\arraybackslash}p{0.58\textwidth}}
\caption{The cited facts Lean reports for each numbered statement, by section (S1 and S2 are
the supplement's).}
\label{tab:statements}\\
\toprule
Section & Cited facts & Statements \\
\midrule
\endfirsthead
\toprule
Section & Cited facts & Statements \\
\midrule
\endhead
\bottomrule
\endlastfoot
% Generated by scripts/theory_axioms_table.py from `#print axioms`; do not edit by hand.
% Printed statements with no \leanok declaration: thm:cited-A30, thm:cited-A13, thm:cited-A15, thm:cited-A19, thm:cited-A20, thm:cited-A55, thm:cited-A57, prop:expanding-grid-pyramid-not-isotropic, prop:grid-pyramid-singular-kernels, prop:band-limited-refinable-mask.
2 & none & Definitions~\ref{def:cascade-measurement-family}, \ref{def:self-decomposable}, \ref{def:isotropic-cascade-family}, \ref{def:isotropic-origin-data}, \ref{def:complete-bernstein}, \ref{def:thorin-bernstein}; Theorems~\ref{thm:cited-A1}, \ref{thm:cited-A34}, \ref{thm:cited-A42}, \ref{thm:cited-A50}, \ref{thm:cited-A51} \\
2 & Theorem~\ref{thm:cited-A22} & Theorems~\ref{thm:cited-A22}, \ref{thm:cited-A23} \\
2 & Theorem~\ref{thm:cited-A24} & Theorem~\ref{thm:cited-A24} \\
2 & Theorem~\ref{thm:cited-A39} & Theorems~\ref{thm:cited-A9}, \ref{thm:cited-A4}, \ref{thm:cited-A31}, \ref{thm:cited-A39}, \ref{thm:cited-A40}, \ref{thm:cited-A44} \\
2 & Theorem~\ref{thm:cited-A41} & Theorem~\ref{thm:cited-A41} \\
2 & Theorem~\ref{thm:cited-A45} & Theorem~\ref{thm:cited-A45} \\
2 & Theorem~\ref{thm:cited-A46} & Theorem~\ref{thm:cited-A46} \\
2 & Theorem~\ref{thm:cited-A48} & Theorem~\ref{thm:cited-A48} \\
2 & Theorem~\ref{thm:cited-A49} & Theorem~\ref{thm:cited-A49} \\
2 & Theorems~\ref{thm:cited-A39}, \ref{thm:cited-A43} & Theorem~\ref{thm:cited-A43} \\
2 & Theorems~\ref{thm:cited-A39}, \ref{thm:cited-A43}, \ref{thm:cited-A45} & Theorem~\ref{thm:cited-A35} \\
3 & none & Definitions~\ref{def:family-dimension-d}, \ref{def:knots-and-ladder}, \ref{def:strata}; Lemmas~\ref{lem:isotropic-symmetric}, \ref{lem:reduction-keeps-stratum}; Propositions~\ref{prop:ladder-not-every-frequency}, \ref{prop:equivalent-weights} \\
3 & Theorem~\ref{thm:cited-A39} & Lemmas~\ref{lem:symmetric-levy-khintchine}, \ref{lem:synthesis}; Proposition~\ref{prop:what-does-not-lift} \\
3 & Theorem~\ref{thm:cited-A41} & Lemma~\ref{lem:steps-infinitely-divisible} \\
3 & Theorem~\ref{thm:cited-A48} & Lemma~\ref{lem:lattice-infinitely-divisible} \\
3 & Theorems~\ref{thm:cited-A39}, \ref{thm:cited-A41} & Lemmas~\ref{lem:family-structure}, \ref{lem:family-dilations}, \ref{lem:discrete-covariance-levy}; Propositions~\ref{prop:stratum-one-example}, \ref{prop:no-lattice-kernel}; Theorem~\ref{thm:lattice-strata} \\
3 & Theorems~\ref{thm:cited-A39}, \ref{thm:cited-A43} & Proposition~\ref{prop:what-does-not-lift-chain} \\
3 & Theorems~\ref{thm:cited-A41}, \ref{thm:cited-A46} & Proposition~\ref{prop:no-bounded-step} \\
3 & Theorems~\ref{thm:cited-A39}, \ref{thm:cited-A41}, \ref{thm:cited-A45} & Proposition~\ref{prop:stratum-two-example}; Theorem~\ref{thm:lattice-strata-strict} \\
3 & Theorems~\ref{thm:cited-A39}, \ref{thm:cited-A41}, \ref{thm:cited-A43}, \ref{thm:cited-A45}, \ref{thm:cited-A48}, \ref{thm:cited-A49} & Proposition~\ref{prop:ladder-knot-kernel} \\
4 & none & Definitions~\ref{def:q-covariant-line}, \ref{def:folded-data}, \ref{def:grid-pyramid-line}; Lemma~\ref{lem:line-axioms-kernels}; Propositions~\ref{prop:recursion-helmholtz}, \ref{prop:sampled-laplace-exact} \\
4 & Theorems~\ref{thm:cited-A39}, \ref{thm:cited-A41}, \ref{thm:cited-A45}, \ref{thm:cited-A46}, \ref{thm:cited-A48}, \ref{thm:cited-A49} & Proposition~\ref{prop:grid-pyramids-infinitely-divisible} \\
5 & none & Definitions~\ref{def:matern-pyramid-data}, \ref{def:frame-matern-pyramid}, \ref{def:implicit-euler-principle}; Lemmas~\ref{lem:bessel-profile}, \ref{lem:bessel-k}, \ref{lem:cm-log-convex}, \ref{lem:cm-exponential-rate}, \ref{lem:cm-uncancelled-atom}, \ref{lem:no-first-order-factorization}; Propositions~\ref{prop:thorin-in-dimension-d}, \ref{prop:matern-generator-octave}, \ref{prop:creation-plane}, \ref{prop:stage-lemma}, \ref{prop:principle-other-readings}, \ref{prop:principle-after-decimation} \\
5 & Theorem~\ref{thm:cited-A39} & Propositions~\ref{prop:matern-increments}, \ref{prop:frame-matern-pyramid-exists} \\
5 & Theorems~\ref{thm:cited-A39}, \ref{thm:cited-A48} & Proposition~\ref{prop:creation-grid} \\
5 & Theorems~\ref{thm:cited-A39}, \ref{thm:cited-A43}, \ref{thm:cited-A45} & Proposition~\ref{prop:laplacian-pyramid} \\
5 & Theorems~\ref{thm:cited-A39}, \ref{thm:cited-A41}, \ref{thm:cited-A43}, \ref{thm:cited-A45} & Theorem~\ref{thm:selection-stratum-two} \\
5 & Theorems~\ref{thm:cited-A22}, \ref{thm:cited-A39}, \ref{thm:cited-A43}, \ref{thm:cited-A45}, \ref{thm:cited-A24} & Theorem~\ref{thm:matern-pyramid} \\
6 & none & Definitions~\ref{def:grid-family}, \ref{def:exact-senses}, \ref{def:exact-grid-pyramid}, \ref{def:level-records}, \ref{def:nine-point-laplacian}; Lemmas~\ref{lem:records-creation}, \ref{lem:cosine-bounds}, \ref{lem:nine-point-symbol}; Propositions~\ref{prop:grid-frame-step-not-selected}, \ref{prop:binomial-records}, \ref{prop:no-isotropic-refinable-finite-mask}, \ref{prop:discrete-helmholtz}, \ref{prop:nine-point-family} \\
6 & Theorems~\ref{thm:cited-A39}, \ref{thm:cited-A41} & Lemma~\ref{lem:knot-data} \\
6 & Theorems~\ref{thm:cited-A41}, \ref{thm:cited-A46} & Proposition~\ref{prop:binomial-knot-relation} \\
6 & Theorems~\ref{thm:cited-A39}, \ref{thm:cited-A41}, \ref{thm:cited-A48} & Propositions~\ref{prop:grid-pyramid-not-isotropic}, \ref{prop:resolvent-masks-exact} \\
6 & Theorems~\ref{thm:cited-A39}, \ref{thm:cited-A41}, \ref{thm:cited-A46}, \ref{thm:cited-A48} & Proposition~\ref{prop:creation-grid-masks} \\
6 & Theorems~\ref{thm:cited-A39}, \ref{thm:cited-A41}, \ref{thm:cited-A48}, \ref{thm:cited-A49} & Proposition~\ref{prop:line-and-arrays} \\
6 & Theorems~\ref{thm:cited-A39}, \ref{thm:cited-A41}, \ref{thm:cited-A46}, \ref{thm:cited-A48}, \ref{thm:cited-A49} & Proposition~\ref{prop:exact-masks} \\
S1 & Theorem~\ref{thm:cited-A39} & Lemma~\ref{lem:line-exponent-data} \\
S1 & Theorems~\ref{thm:cited-A39}, \ref{thm:cited-A41} & Lemmas~\ref{lem:closed-subgroups}, \ref{lem:discrete-covariance-line}; Propositions~\ref{prop:lindeberg-stratum-one}, \ref{prop:discrete-gaussian-pyramid}, \ref{prop:no-lattice-kernel-line}; Theorems~\ref{thm:lattice-line}, \ref{cor:lattice-line-strict} \\
S2 & none & Definitions~\ref{def:modulation-orbit}, \ref{def:phase-entropy}; Lemmas~\ref{lem:gauge-group}, \ref{lem:averaging-bound}, \ref{lem:null-conull-point-mass}, \ref{lem:kl-measurable}, \ref{lem:phase-disintegration-exists}, \ref{lem:haar-smearing-density}, \ref{lem:haar-smearing-injective}, \ref{lem:laplace-atoms-admissible}; Propositions~\ref{prop:phase-entropy-scale-unit}, \ref{prop:classical-pyramids-approximate}, \ref{prop:extreme-rays-random-walks}; Theorem~\ref{thm:knot-atom-cascades} \\
S2 & Theorem~\ref{thm:cited-A39} & Proposition~\ref{prop:strata-cylinder-coordinates} \\
S2 & Theorems~\ref{thm:cited-A22}, \ref{thm:cited-A24} & Theorem~\ref{thm:matern-pyramid-line-noncreation} \\
S2 & Theorems~\ref{thm:cited-A22}, \ref{thm:cited-A39} & Proposition~\ref{prop:ladder-non-creation-line} \\
S2 & Theorems~\ref{thm:cited-A39}, \ref{thm:cited-A41} & Lemmas~\ref{lem:step-field-determines-family}, \ref{lem:knot-exponent-ladder}, \ref{lem:member-gaussian-part}; Propositions~\ref{prop:modulated-selection-strict}, \ref{prop:appearance-entropy}; Theorems~\ref{thm:annulus-parametrization}, \ref{thm:annulus-parametrization-phase}, \ref{thm:surface-measure}, \ref{thm:phase-entropy-selection}, \ref{thm:phase-entropy-modulated-selection}, \ref{thm:matern-pyramid-line} \\
S2 & Theorems~\ref{thm:cited-A39}, \ref{thm:cited-A43} & Proposition~\ref{prop:strata-cylinder-coordinates-q-cin} \\
S2 & Theorems~\ref{thm:cited-A41}, \ref{thm:cited-A46} & Proposition~\ref{prop:classical-kernels-no-exact-cascade} \\
S2 & Theorems~\ref{thm:cited-A39}, \ref{thm:cited-A41}, \ref{thm:cited-A24} & Proposition~\ref{prop:classical-pyramids-approximate-discrete-gaussian} \\
S2 & Theorems~\ref{thm:cited-A39}, \ref{thm:cited-A41}, \ref{thm:cited-A48} & Proposition~\ref{prop:phase-entropy-ladder-bound} \\

\end{longtable}
}
